% Terjemahan appendix2.tex, baris sumber 606--660.
\section{Teorema Pembenaman Freyd--Mitchell}\label{sec:FM}
Kita tetap memilih suatu semesta Grothendieck $\mathcal{U}$ untuk membedakan kategori (yang secara bawaan berarti kategori-$\mathcal{U}$) dari kategori kecil. Sesuai konvensi, semua grup, gelanggang, dan modul yang disebutkan di bawah ini secara bawaan direalisasikan pada himpunan kecil (yakni himpunan-$\mathcal{U}$).

\begin{lemma}\label{prop:FM-prep-1}
	Misalkan $\mathcal{C}$ suatu kategori abelian kokomplet yang memiliki pembangkit projektif, dan $O$ suatu himpunan bagian kecil dari $\Obj(\mathcal{C})$. Maka terdapat pembangkit projektif $S$ dari $\mathcal{C}$ sedemikian sehingga setiap $X \in O$ dapat dinyatakan sebagai objek hasil bagi dari $S$.
\end{lemma}
\begin{proof}
	Ambil sebarang pembangkit projektif $s$ dari $\mathcal{C}$. Karena $\mathcal{C}$ kokomplet, bentuk jumlah langsung
	\[ S := \bigoplus_{X \in O} s^{\oplus \Hom(s, X)}. \]
	Objek ini tetap projektif. Untuk setiap $X \in O$, definisikan morfisme kanonik $S \twoheadrightarrow X$ yang restriksinya pada suku jumlah langsung $s$ yang bersesuaian dengan $f \in \Hom(s, X)$ sama dengan $f$, dan pada semua suku jumlah langsung lainnya sama dengan $0$. Dari sifat pembangkit mudah dilihat bahwa ini merupakan epimorfisme; lihat pembuktian Teorema \ref{prop:identify-ModR} untuk rinciannya.
\end{proof}

Hasil berikut merupakan suatu variasi sederhana dari Teorema \ref{prop:identify-ModR}.

\begin{lemma}\label{prop:FM-prep-2}
	Misalkan $S$ pembangkit projektif dari kategori abelian $\mathcal{C}$. Definisikan $R := \End_{\mathcal{C}}(S)$, sehingga diperoleh funktor setia dan eksak (Proposisi \ref{prop:injective-cogenerator})
	\[ G := \Hom_{\mathcal{C}}(S, \cdot): \mathcal{C} \to \cated{Mod}R. \]
	Jika $X \in \Obj(\mathcal{C})$ dapat dinyatakan sebagai hasil bagi dari jumlah langsung sejumlah hingga salinan $S$, maka pemetaan yang diinduksi oleh $G$
	\[ \Hom_{\mathcal{C}}(X, Y) \to \Hom_{\cated{Mod}R}(GX, GY) \]
	bijektif untuk setiap $Y \in \Obj(\mathcal{C})$.
\end{lemma}
\begin{proof}
	Ambil $m \in \Z_{\geq 0}$ dan barisan eksak pendek dalam $\mathcal{C}$
	\[ 0 \to X' \to S^{\oplus m} \to X \to 0. \]
	Menerapkan $G$ padanya memberikan barisan eksak pendek dalam $\cated{Mod}R$. Tinjau diagram komutatif dengan baris-baris eksak dalam $\cate{Ab}$
	\begin{equation}\label{eqn:FM-prep-2}\begin{tikzcd}
			0 \arrow[r] & \Hom_{\mathcal{C}}(X, Y) \arrow[r] \arrow[hookrightarrow, d] & \Hom_{\mathcal{C}}(S^{\oplus m}, Y) \arrow[r] \arrow[hookrightarrow, d] & \Hom_{\mathcal{C}}(X', Y) \arrow[hookrightarrow, d] \\
			0 \arrow[r] & \Hom_R(GX, GY) \arrow[r] & \Hom_R(G(S^{\oplus m}), GY) \arrow[r] & \Hom_R(GX', GY),
	\end{tikzcd}\end{equation}
	semua panah vertikal berasal dari funktor setia $G$, sehingga merupakan injeksi. Perhatikan bahwa komposisi
	\[ \Hom_{\mathcal{C}}(S, Y) \xrightarrow{\identity} GY \xleftarrow[\psi(1_R) \mapsfrom \psi]{\sim} \Hom_R(R, GY) = \Hom_R(GS, GY) \]
	sama dengan homomorfisme pada $\Hom$ yang diinduksi oleh funktor $G$. Tentu saja ini nyaris suatu tautologi, dan telah diperiksa dalam pembuktian Lema \ref{prop:GP-prep}. Karena $G$ merupakan funktor aditif, panah vertikal di bagian tengah \eqref{eqn:FM-prep-2} adalah isomorfisme. Dengan mengejar diagram dalam $\cate{Ab}$, mudah dilihat bahwa panah vertikal di sebelah kiri juga merupakan isomorfisme.
\end{proof}

\begin{theorem}[P.\ Freyd, B.\ Mitchell {\cite[Bab 4]{Fr03}}]\label{prop:FM}
	\index{Freyd-Mitchell@Teorema Pembenaman Freyd--Mitchell (Freyd--Mitchell Embedding Theorem)}
	Misalkan $\mathcal{A}$ kategori abelian kecil. Maka terdapat gelanggang $R$ dan funktor penuh, setia, dan eksak $F: \mathcal{A} \to \cated{Mod}R$.
\end{theorem}
\begin{proof}
	Argumen berikut diambil dari \cite[Teorema 9.6.10]{KS06}. Karena $\mathcal{A}^{\opp}$ tetap merupakan kategori abelian kecil, Teorema \ref{prop:Ind-Grothendieck} menunjukkan bahwa $\IndC(\mathcal{A}^{\opp})$ adalah kategori Grothendieck. Dengan demikian, Korolari \ref{prop:Grothendieck-cat-cogenerator} menjamin bahwa kategori tersebut memiliki kogenerator injektif. Mengingat
	\[ \ProC(\mathcal{A}) \simeq \IndC(\mathcal{A}^{\opp})^{\opp}, \]
	kita menyimpulkan bahwa $\ProC(\mathcal{A})$ memiliki pembangkit projektif. Selain itu, $\mathcal{A} \to \ProC(\mathcal{A})$ merupakan funktor penuh, setia, dan eksak di antara kategori-kategori abelian (Teorema \ref{prop:Ind-Abel}).
	
	Pandang $\mathcal{A}$ sebagai subkategori penuh dari $\ProC(\mathcal{A})$. Dengan menerapkan Lema \ref{prop:FM-prep-1} (ambil $O = \Obj(\mathcal{A})$), diperoleh pembangkit projektif $S$ dari $\ProC(\mathcal{A})$ sedemikian sehingga setiap objek dari $\mathcal{A}$ dapat dinyatakan sebagai hasil bagi dari $S$. Sekarang tinjau funktor
	\[ \mathcal{A} \to \ProC(\mathcal{A}) \xrightarrow{G} \cated{Mod}R , \quad G := \Hom_{\ProC(\mathcal{A})}(S, \cdot), \; R := \End_{\ProC(\mathcal{A})}(S). \]\footnote{Catatan penerjemah: sumber mencetak subskrip gelanggang endomorfisme sebagai $\ProC\mathcal{C}$. Dalam konteks ini, $S$ adalah objek $\ProC(\mathcal{A})$, sehingga subskrip yang bertipe benar ialah $\ProC(\mathcal{A})$.}
	Tuliskan komposisinya sebagai $F$. Karena setiap ruas bersifat eksak, $F$ merupakan funktor eksak. Dengan menerapkan Lema \ref{prop:FM-prep-2}, diketahui bahwa $G$ penuh dan setia; jadi, $F$ penuh dan setia. Ini membuktikan pernyataan tersebut.
\end{proof}

\begin{corollary}
	Misalkan $\mathcal{A}$ kategori abelian kecil. Maka terdapat funktor setia dan eksak $E: \mathcal{A} \to \cate{Ab}$.
\end{corollary}

Teorema Freyd--Mitchell memungkinkan banyak diagram komutatif mengenai kategori abelian umum direduksi untuk diperiksa dalam kasus $\cated{Mod}R$, atau bahkan dalam $\cate{Ab}$. Karena kedua kategori terakhir bersifat konkret, objek-objeknya memiliki elemen, dan teknik pengejaran diagram \cite[\S 6.8]{Li1} sangat menyederhanakan banyak persoalan. Di bawah asumsi teori himpunan yang sesuai, prasyarat bahwa kategorinya kecil dapat dipenuhi dengan memperbesar semesta Grothendieck.
\index{persoalan ukuran}
